\documentclass[12pt]{article}
\usepackage[utf8]{inputenc}
\usepackage[T1]{fontenc}
\usepackage{hyperref}
\title{Zapotec, Santa Catarina Albarradas Language Resources for K-12 Curriculum Development}
\author{Compiled by Claude Sonnet 5.5 (claude-sonnet-5-5) for Andr\'as Kornai}
\date{}
\begin{document}
\maketitle

\section{Introduction}
Santa Catarina Albarradas Zapotec (ISO 639-3 ztn, Glottocode sant1448) is an Otomanguean language of Oaxaca, Mexico, also called San Antonio Albarradas Zapotec. It is in Central Core Zapotec. A 1990 census figure gives about 1,000 speakers. Joshua Project reports about 1,600 people in the group. Wikitongues rates it EGIDS 6b (Threatened). The searches found few resources.

\section{Primary Educational Resources}
No primers, textbooks or school programmes were found.

\section{Linguistic and Reference Materials}
\begin{itemize}
\item Wikipedia overview (notes partial intelligibility with Santo Domingo Albarradas Zapotec): \url{https://en.wikipedia.org/wiki/Santa_Catarina_Albarradas_Zapotec}
\item Glottolog entry: \url{https://glottolog.org/resource/languoid/id/sant1448.bigmap.html}
\item Wikitongues language page: \url{https://wikitongues.org/languages/ztn/}
\end{itemize}
No dictionary or grammar was found.

\section{Religious and Community Resources}
\begin{itemize}
\item Joshua Project people-group profile: \url{https://joshuaproject.net/people_groups/14735} (a missions organisation, so its framing is religious)
\end{itemize}

\section{Digital and Audio Resources}
\begin{itemize}
\item Global Recordings Network audio, including a Bible-story series: \url{https://www.globalrecordings.net/en/language/ztn}
\end{itemize}

\section{Curriculum Development Context}
The language is still used by adults, which gives a real base for work. The documentation found is thin, so curriculum development would need an orthography, a lexicon and recorded texts developed with the community. The neighbouring Albarradas variety may be a partial model, but it is a different language.
\end{document}